\chapter{52}



Gegen acht würde es eindunkeln. Bis dahin blieb Louis ausreichend Zeit, sich
umzusehen. 

Mit dem Küchentuch über der Schulter trat er vor die Hütte. Er liess
seinen Blick den Bergen entlang schweifen und wunderte sich von Neuem.
Über die Schönheit des Ortes. Über die Selbstverständlichkeit, wie Joh
und Manuela ihm das kleine Haus überliessen. Über sein Glück, hier zu
sein. Ein unwirklicher Friede schien es sich in jedem noch so kleinen
Winkel hier oben bequem gemacht zu haben.
Louis atmete tief durch und betrat die Hütte. Er spülte das Geschirr.
Nach dem Wegräumen des letzten Glases war es soweit. Bis auf seine kurze
Hose zog er sich aus, schnappte sich Handy und Kopfhörer und trat
barfuss hinaus auf den kargen, trockenen Platz neben
der Hütte.

«Wem nützt es, wenn ich die Stimmung nicht geniesse. Nichts war, nichts
kommt. Alles ist gut.»

Daran hängte er ein bestimmtes \emph{ja.} Er rief es laut ins Tal
hinunter und fing augenblicklich an zu lachen. Er redete wieder
mit sich selbst.
Louis suchte sich durch die Titel von «Álvaro Soler». Bei
\emph{«Volar»} blieb er stecken und startete den Song. Er hatte
Herzklopfen.

Die ersten Takte nahmen ihn direkt an die Hand. Seine Beine begannen den
Rhythmus zu tanzen. Erst zögernd. Dann hielt er sich
nicht mehr zurück. Louis segelte, seine Arme ausgestreckt, ähnlich eines
Flugzeuges, wie er es als kleines Kind imitiert hatte, rund um die Hütte
und anschliessend in alle Richtungen. Dazu sang er den Text lauthals
mit. Niemand konnte ihn sehen. Dass er den Song einige Male seinen
Leuten vorgesungen hatte, spielte gerade keine Rolle. Ebenso wenig, dass
Giorgia ihn mochte. Gerade gehörte er nur ihm.

\qrblock{Álvaro Soler \emph{«Volar»}}{media/QR-Codes/035_Volar_Àlvaro_Soler_Spotify.png}

Es musste an dieser hochdosierten Einfachheit liegen. Louis’ lange
verlorener Übermut kam zurück. «Ich will mehr, ich will mehr. Und
hab’ keine Minute zu verlieren», ja, so konnte man es
auch sehen. All die üblen Situationen und damit verbundenen Gefühle mussten doch über kurz oder lang abzulegen sein. Andere schafften das doch auch.
Nachdem er stehen geblieben und den nächsten Song angetippt hatte, ging
es in gleicher Manier weiter. \emph{«Sofia»}  sog ihn auf wie einen kleinen Fisch und wirbelte ihn in alle Richtungen. Schliesslich sass er erfüllt am Boden. Wippte mit den Füssen den Takt und konzentrierte sich auf den Text. Er war verblüfft.
So bewusst hatte er bisher nicht hingehört. Soler besang das Ende einer Liebe. Mit all den herausfordernden Facetten des Abschieds. Doch die
schmissige Musik trug die Beschreibung mit Optimismus. Louis liess sich
nach hinten fallen. Schliesslich lag er flach auf der Wiese. Er streckte
die Arme weit von sich und startete «Sofia» erneut.
Es gab so viel Platz hier. Mit geschlossenen Augen summte er die Melodie
mit.

\qrblock{Álvaro Soler \emph{«Sofia»}}{media/QR-Codes/036_Sofia_Àlvaro_Soler_Spotify.png}

Noch während des dritten Durchlaufs sank er in entspanntes Dösen und
schlief schliesslich ein.

\sectionbreak

Mit einem Schaudern wachte er auf. Ihn fröstelte. Sein Rücken schmerzte
und für einen Moment fragte er sich, wo er war.
Das Licht war anders. Intimer. Als habe jemand über dem Ort die Vorhänge zugezogen. 
Louis zupfte die Kopfhörer aus den Ohren und blickte auf sein Handy. Er
wunderte sich. Erschöpfung und Ruhe fernab schienen eine ideale
Kombination zu sein, sich fallenzulassen.
Er fühlte sich gerädert, war müder als vorher, stand schwerfällig auf
und steuerte auf die Hütte zu. Gänsehaut überzog seinen Oberkörper. Auf
der Bank vor dem Eingang schnappte er sich hastig sein T-Shirt. 
